\documentclass[12pt]{article}
\usepackage[utf8]{inputenc}
\usepackage{amsmath, amssymb, amsthm}
\usepackage{geometry}
\usepackage{fancyhdr}

\geometry{letterpaper, margin=1in}
\setlength{\headheight}{15pt}

% Header setup
\pagestyle{fancy}
\fancyhf{}
\lhead{Math 311 Fall 2026}
\rhead{Jeffrey Lu}
\cfoot{\thepage}

\title{Homework Assignment 3 \\ \large Math 311, Fall 2026 \\ Rutgers University}
\author{}
\date{}

\begin{document}

\maketitle

\noindent{} \textbf{Exercise 2.2.7}
\begin{enumerate}
    \item [(a)] $(-1)^n$ is frequently in the set $\{1\}$ since any even $n$ gives $(-1)^n = 1$. But \(((-1)^n)\) is not eventually in \(\{1\}\), since every odd-indexed term equals \(-1\).
    \item [(b)] Eventually is stronger since eventually $ \implies $ frequently but not the converse.
    \item [(c)] We can say \((a_n)\) converges to \(a\) if,
          for every \(\varepsilon>0\), the sequence is eventually in \(V_\varepsilon(a)\).
    \item [(d)] Yes, $(x_n)$ is frequently in $(1.9,2.1)$ because infinitely many terms are equal to 2 in the interval.
          However, it might not be eventually. For example, consider
          \[
              x_n =
              \begin{cases}
                  2, & n \text{ even}, \\
                  0, & n \text{ odd}.
              \end{cases}
          \]
\end{enumerate}

\noindent{} \textbf{Exercise 2.3.2}
\begin{enumerate}
    \item [(a)]
          Let \(\varepsilon>0\). Choose \(N\) such that, for \(n\ge N\),
          \[
              |x_n-2|<\frac{3\varepsilon}{2}.
          \]Then, for \(n\ge N\),
          \[
              \left|\frac{2x_n-1}{3}-1\right|
              =\left|\frac{2x_n-4}{3}\right|
              =\frac{2}{3}|x_n-2|
              <\frac{2}{3}\cdot\frac{3\varepsilon}{2}
              =\varepsilon.
          \]Therefore,
          \[
              \frac{2x_n-1}{3}\to 1.
          \]
    \item [(b)]
          Let \(\varepsilon>0\). Because \(x_n\to2\), choose \(N\) such that for \(n\ge N\),
          \[
              |x_n-2|<\min\{1,2\varepsilon\}.
          \]The condition \(|x_n-2|<1\) implies \(x_n>1\), so \(|x_n|>1\). Thus, for \(n\ge N\),
          \[
              \left|\frac1{x_n}-\frac12\right|
              =
              \left|\frac{2-x_n}{2x_n}\right|
              =
              \frac{|x_n-2|}{2|x_n|}
              <
              \frac{|x_n-2|}{2}.
          \]Since \(|x_n-2|<2\varepsilon\),
          \[
              \left|\frac1{x_n}-\frac12\right|
              <
              \frac{2\varepsilon}{2}
              =
              \varepsilon.
          \]Hence,
          \[
              \frac1{x_n}\to\frac12.
          \]
\end{enumerate}

\noindent{} \textbf{Exercise 2.3.5}
Let
\[
    (z_n)=(x_1,y_1,x_2,y_2,x_3,y_3,\ldots).
\]
Suppose $(z_n) \to L$. Let $\varepsilon > 0$. There exists $N$ such that if $m \geq N$, then
\[
    |z_m - L| < \varepsilon.
\]
For $n \geq N$, we have $2n-1 \geq N$, so
\[
    |x_n - L| = |z_{2n-1} - L| < \varepsilon.
\]
Thus, $x_n \to L.$ Similarly, if \(n\ge N\), then \(2n\ge N\), and hence
\[
    |y_n-L|=|z_{2n}-L|<\varepsilon.
\]
So \(y_n\to L\).
Therefore, both \((x_n)\) and \((y_n)\) converge and have the same limit.

Conversely, suppose
\[
    x_n\to L
    \qquad\text{and}\qquad
    y_n\to L.
\]
Let $\varepsilon > 0$. Choose $N_1$ such that if $n \geq N_1$, then $|x_n - L| < \varepsilon$,
and choose $N_2$ such that if $n \geq N_2$, then $|y_n - L| < \varepsilon$.
Let $N = 2 \cdot \max\{N_1, N_2\}$. Then for any $m \geq N$, we have even/odd cases.

If $m = 2k$, then $k \geq \max\{N_1, N_2\}$, so
\[
    |z_m-L|=|z_{2k}-L|=|y_k-L|<\varepsilon.
\]

If $m = 2k-1$, then $k \geq \max\{N_1, N_2\}$, so
\[
    |z_m-L|=|z_{2k-1}-L|=|x_k-L|<\varepsilon.
\]
Hence \(z_m\to L\).


\noindent{} \textbf{Exercise 2.3.6}
Notice
\[
    b_n=n-\sqrt{n^2+2n}
    =\frac{n^2-(n^2+2n)}{n+\sqrt{n^2+2n}}
    =\frac{-2n}{n+\sqrt{n^2+2n}}.
\]
Because \(n>0\), factor \(n\) out of the square root:
\[
    \sqrt{n^2+2n}
    =
    \sqrt{n^2\left(1+\frac{2}{n}\right)}
    =
    n\sqrt{1+\frac{2}{n}}.
\]Therefore,
\[
    b_n
    =
    \frac{-2n}{n+n\sqrt{1+\frac{2}{n}}}
    =
    \frac{-2}{1+\sqrt{1+\frac{2}{n}}}.
\]Now,
\[
    \frac1n\to0,
\]so by the Algebraic Limit Theorem,
\[
    1+\frac2n\to1.
\]
and
\[
    \sqrt{1+\frac2n}\to\sqrt1=1.
\]
Thus the denominator converges:
\[
    1+\sqrt{1+\frac2n}\to 1+1=2.
\]Using the quotient rule from the Algebraic Limit Theorem,
\[
    \lim_{n\to\infty} b_n
    =
    \frac{-2}{2}
    =
    -1.
\]
Hence \(\lim b_n\) exists and equals -1.

\noindent{}\textbf{Exercise 2.3.8}
\begin{enumerate}
    \item [(a)]
          Since \(x_n\to x\), the Algebraic Limit Theorem gives
          \[
              x_n^k \to x^k
              \qquad (k=1,2,\ldots,d).
          \]
          Hence
          \[
              a_kx_n^k\to a_kx^k
          \]
          for every \(k\).
          Therefore,
          \[
              p(x_n)
              =
              a_0+a_1x_n+\cdots+a_dx_n^d
              \to
              a_0+a_1x+\cdots+a_dx^d
              =
              p(x).
          \]
    \item [(b)]
          Let
          \[
              f(t)=
              \begin{cases}
                  0, & t=0,    \\
                  1, & t\ne 0,
              \end{cases}
          \]and choose
          \[
              x_n=\frac1n.
          \]Then \(x_n\to0\), but \(x_n\ne0\) for every \(n\). Hence
          \[
              f(x_n)=1
          \]for all \(n\), so \((f(x_n))\to1\). However,
          \[
              f(0)=0.
          \]Thus \(f(x_n)\) converges, but not to \(f(x)\).
\end{enumerate}

\noindent{} \textbf{Exercise 2.3.11}
\begin{enumerate}
    \item [(a)]Suppose \(x_n\to x\), and let
          \[
              y_n=\frac{x_1+x_2+\cdots+x_n}{n}.
          \]Then
          \[
              y_n-x
              =
              \frac{(x_1-x)+(x_2-x)+\cdots+(x_n-x)}{n}.
          \]
          Let $\varepsilon > 0$.
          Since $x_n \to x$, there exists $N$ such that if $n \geq N$ then
          \[
              |x_n - x| < \frac{\varepsilon}{2}.
          \]
          Define
          \[
              C=\sum_{k=1}^{N-1}|x_k-x|.
          \]

          For $n \geq N$, we have
          \[
              \begin{aligned}
                  |y_n - x| & \leq \frac{1}{n}\sum_{k=1}^{n}|x_k - x|             \\
                            & = \frac{C}{n} + \frac{1}{n}\sum_{k=N}^{n}|x_k - x|.
              \end{aligned}
          \]
          Every term in the second sum is less than \(\varepsilon/2\), so
          \[
              \frac{1}{n}\sum_{k=N}^n |x_k-x|
              \leq \frac{n-N+1}{n}\cdot\frac{\varepsilon}{2}
              \leq \frac{\varepsilon}{2}.
          \]
          Choose \(n\) large enough that \(C/n<\varepsilon/2\). Then
          \[
              |y_n-x|<\frac{\varepsilon}{2}+\frac{\varepsilon}{2}
              =\varepsilon.
          \]
          Therefore \(y_n\to x\).
    \item [(b)]
          Let
          \[
              x_n=(-1)^n.
          \]This sequence does not converge. But its averages do:
          \[
              y_n=\frac{(-1)^1+(-1)^2+\cdots+(-1)^n}{n}.
          \]If \(n\) is even, the numerator is \(0\), so \(y_n=0\). If \(n\) is odd, the numerator is \(-1\), so
          \[
              y_n=-\frac1n.
          \]Thus \(y_n\to0\), even though \((x_n)\) does not converge.
\end{enumerate}

\noindent{}\textbf{Exercise 2.4.1}

Let
\[
    \alpha=2-\sqrt3.
\]Notice that \(\alpha\) satisfies
\[
    \alpha=\frac1{4-\alpha},
\]because \(\alpha^2-4\alpha+1=0\).
\begin{enumerate}
    \item [(a)]
          First,
          \[
              x_2=\frac1{4-x_1}=\frac1{4-3}=1.
          \]
          Also, \(\alpha<1=x_2\). Suppose \(x_n>\alpha\) and \(x_n\le 1\). Then
          \[
              x_{n+1}-\alpha
              =
              \frac1{4-x_n}-\frac1{4-\alpha}
              =
              \frac{x_n-\alpha}{(4-x_n)(4-\alpha)}>0.
          \]
          So \(x_{n+1}>\alpha\). Also, if \(x_n\le1\), then
          \[
              x_{n+1}=\frac{1}{4-x_n}\le\frac13<1.
          \]
          Hence, by induction, \(\alpha<x_n\le1\) for every \(n\ge2\).
          Now, for \(n\ge2\),
          \[
              x_{n+1}-x_n
              =
              \frac1{4-x_n}-x_n
              =
              \frac{x_n^2-4x_n+1}{4-x_n}.
          \]Factor the numerator:
          \[
              x_n^2-4x_n+1
              =
              (x_n-(2-\sqrt3))(x_n-(2+\sqrt3)).
          \]Because
          \[
              2-\sqrt3<x_n\le1<2+\sqrt3,
          \]the numerator is negative, while \(4-x_n>0\). Hence
          \[
              x_{n+1}-x_n<0.
          \]Therefore \((x_n)\) is decreasing and bounded below by \(2-\sqrt3\). By the Monotone Convergence Theorem, \((x_n)\) converges.
    \item [(b)] If \(x_n\to L\), then removing the first term does not change the limit.
    \item [(c)] Taking limits in
          \[
              x_{n+1}=\frac1{4-x_n},
          \]
          gives
          \[
              L=\frac1{4-L}.
          \]
          Thus
          \[
              L(4-L)=1,
          \]
          so
          \[
              L^2-4L+1=0.
          \]
          The possible solutions are
          \[
              L=2\pm\sqrt3.
          \]
          But \(x_n\le1\) for all \(n\ge2\), so its limit must satisfy \(L\le1\). Therefore,
          \[
              \lim_{n\to\infty}x_n=2-\sqrt3.
          \]
\end{enumerate}

\noindent{} \textbf{Exercise 2.4.2}
\begin{enumerate}
    \item [(a)] The problem is that the argument assumes the limit exists.
          In fact,
          \[
              y_1=1,\qquad y_2=3-1=2,\qquad y_3=3-2=1,
          \]so the sequence is
          \[
              1,2,1,2,1,2,\ldots
          \]and does not converge.
    \item [(b)] Only if we prove convergence is true first.
          Let
          \[
              \beta=\frac{3+\sqrt5}{2}.
          \]
          Notice that
          \[
              y_1=1<\beta.
          \]

          Assume \(1\le y_n<\beta\). Then
          \[
              y_{n+1}=3-\frac1{y_n}\ge 2>1.
          \]
          Also, because \(0<y_n<\beta\),
          \[
              \frac1{y_n}>\frac1\beta,
          \]
          and hence
          \[
              y_{n+1}
              =3-\frac1{y_n}
              <3-\frac1\beta
              =\beta.
          \]
          Thus \(1\le y_n<\beta\) for every \(n\).

          Also,
          \[
              y_{n+1}-y_n
              =
              3-\frac1{y_n}-y_n
              =
              \frac{-(y_n^2-3y_n+1)}{y_n}.
          \]
          Since
          \[
              y_n^2-3y_n+1
              =
              \left(y_n-\frac{3-\sqrt5}{2}\right)
              \left(y_n-\frac{3+\sqrt5}{2}\right),
          \]
          and \(1\le y_n<\beta\), we have \(y_{n+1}-y_n>0\).
          Thus \((y_n)\) is increasing and bounded above, so it converges.

          Let
          \[
              L=\lim_{n\to\infty}y_n.
          \]
          Then
          \[
              L=3-\frac1L,
          \]
          so
          \[
              L^2-3L+1=0.
          \]
          Therefore,
          \[
              L=\frac{3\pm\sqrt5}{2}.
          \]
          Because \(y_n\ge1\), we have
          \[
              \lim_{n\to\infty}y_n=\frac{3+\sqrt5}{2}.
          \]
\end{enumerate}

\noindent{} \textbf{Exercise 2.4.3}
\begin{enumerate}
    \item [(a)] Define
          \[
              x_1=\sqrt2,
              \qquad
              x_{n+1}=\sqrt{2+x_n}.
          \]
          Notice
          \[
              x_1=\sqrt2<2.
          \]If \(0<x_n<2\), then
          \[
              x_{n+1}=\sqrt{2+x_n}<\sqrt4=2.
          \]Thus \(x_n<2\) for every \(n\). Also, when \(0<x_n<2\),
          \[
              2+x_n>x_n^2,
          \]because
          \[
              2+x_n-x_n^2=(2-x_n)(x_n+1)>0.
          \]Taking square roots gives
          \[
              x_{n+1}=\sqrt{2+x_n}>x_n.
          \]Hence \((x_n)\) is increasing and bounded above by \(2\), so it converges.
          Let \(L=\lim x_n\). Then
          \[
              L=\sqrt{2+L}.
          \]Thus
          \[
              L^2-L-2=0,
          \]so
          \[
              L=2 \quad\text{or}\quad L=-1.
          \]Since \(x_n\ge\sqrt2\) for every \(n\), we have \(L\ge\sqrt2>0\).
    \item [(b)]
          Define
          \[
              y_1=\sqrt2,
              \qquad
              y_{n+1}=\sqrt{2y_n}.
          \]Again, \(y_1<2\). If \(0<y_n<2\), then
          \[
              y_{n+1}=\sqrt{2y_n}<\sqrt4=2.
          \]So \(y_n<2\) for every \(n\). Furthermore, if \(0<y_n<2\), then
          \[
              2y_n>y_n^2.
          \]Taking square roots gives
          \[
              y_{n+1}=\sqrt{2y_n}>y_n.
          \]Therefore \((y_n)\) is increasing and bounded above by \(2\), and hence it converges.
          Let \(L=\lim y_n\). Taking limits in
          \[
              y_{n+1}=\sqrt{2y_n}
          \]gives
          \[
              L=\sqrt{2L}.
          \]Thus
          \[
              L^2=2L,
          \]so \(L=0\) or \(L=2\). Since \(y_n\ge\sqrt2>0\),
          \[
              \lim y_n=2.
          \]
\end{enumerate}

\noindent{} \textbf{Exercise 2.4.8}
\begin{enumerate}
    \item [(a)] Let
          \[
              s_N=\sum_{n=1}^N\frac1{2^n}
              =\frac12+\frac14+\cdots+\frac1{2^N}.
          \]This is a geometric sum, so
          \[
              s_N=\frac{\frac12(1-(1/2)^N)}{1-\frac12}
              =1-\frac1{2^N}.
          \]Therefore,
          \[
              \lim_{N\to\infty}s_N=1.
          \]So the series converges, and
          \[
              \sum_{n=1}^{\infty}\frac1{2^n}=1.
          \]
    \item [(b)] Let
          \[
              \begin{aligned}
                  s_N
                   & = \sum_{n=1}^N\left(\frac{1}{n(n+1)}\right)             \\
                   & = \sum_{n=1}^N\left(\frac1n-\frac1{n+1}\right)          \\
                   & = \left(1-\frac12\right) + \left(\frac12-\frac13\right)
                  + \cdots + \left(\frac1N-\frac1{N+1}\right).
              \end{aligned}
          \]
          Everything cancels except the first and last terms:
          \[
              s_N=1-\frac1{N+1}=\frac{N}{N+1}.
          \]Hence,
          \[
              \lim_{N\to\infty}s_N=1.
          \]So
          \[
              \sum_{n=1}^{\infty}\frac1{n(n+1)}=1.
          \]
    \item [(c)]
          Let \[
              s_N
              =
              \sum_{n=1}^N \log\left(\frac{n+1}{n}\right)
              =
              \sum_{n=1}^N\left(\log(n+1)-\log n\right).
          \]
          Then
          \[
              s_N
              =
              (\log2-\log1)+(\log3-\log2)+\cdots+(\log(N+1)-\log N).
          \]
          Thus,
          \[
              s_N=\log(N+1)-\log1=\log(N+1).
          \]
          Since
          \[
              \log(N+1)\to\infty,
          \]the partial sums do not converge to a real number. Therefore the series diverges.
\end{enumerate}
\end{document}